% ============================================================
%  AI 工具使用详情（支撑材料模板）
%  注意：本页不应作为参赛论文正文页面提交。
%  如使用 AI 工具，请将本片段复制到单独的支撑材料说明文档中，
%  并以 PDF 文件“AI 工具使用详情”随支撑材料提交。
% ============================================================
\thispagestyle{plain}

\begin{center}
  {\heiti\zihao{3} AI 工具使用详情}
\end{center}

\vspace{1em}

\noindent 本说明用于记录参赛过程中使用的人工智能工具。未使用 AI 工具时，不需要提交本文件，只需在论文参考文献后保留未使用声明。

\vspace{0.8em}

\section*{一、所用 AI 工具名称和版本}

\begin{itemize}[leftmargin=2em]
  \item 工具名称：[如 ChatGPT / DeepSeek / 通义千问]
  \item 版本或模型：[如 GPT-5 / DeepSeek-R1-0528]
  \item 开发机构或公司：[填写机构名称]
\end{itemize}

\section*{二、具体使用目的和环节}

说明 AI 工具用于哪些环节，例如文献线索整理、代码调试、语言润色、图表说明草拟等。

\section*{三、关键交互记录}

列出重要提示词与回复摘要，尤其是对建模、算法、代码或结论产生影响的交互。

\section*{四、采纳和人工修改情况}

说明哪些内容被采纳，队员如何核查、重写、推导、运行或修改，确保核心建模与分析由参赛队独立完成。
